In this work, we studied \textit{Asymmetric Shapley Values} (ASV), a variant of Shapley values that incorporates causal information into feature attribution. We analyzed the computational complexity of ASV and showed tractability results for restricted settings in which computing SHAP was intractable. To make ASV more practical, we developed exact and approximate algorithms based on equivalence classes of topological orderings. Our experiments show that these techniques improve the performance for computing the score and yields accurate approximations.

Overall, our results suggest that ASV is a promising explainability framework when causal structure is relevant, since it can capture dependencies that standard SHAP does not reflect. Our experiments also indicate that ASV can be implemented effectively for several networks, although its scalability depends strongly on the causal graph. While equivalence classes can greatly reduce the number of topological orderings considered, they may still grow quickly in large or dense DAGs, where the tree or polytree assumptions used by the exact algorithms no longer apply and inference can become more expensive. In these cases, the method may need to rely on sampling-based approximations, whose accuracy depends on both topological-order sampling and conditional-expectation estimation. Thus, further optimization is needed for very large or highly connected graphs.

Several directions remain open for future work. On the theoretical side, it would be interesting to generalize the equivalence class algorithm from directed trees to polytrees, and to further study the complexity of computing the ASV in other cases. On the practical side, possible improvements include implementing alternative counting algorithms from the literature, optimizing the current implementation through caching or parallelization, and extending the framework to more general causal models beyond Bayesian networks.